% META: {"id":"1SPE-EXPO-CO-035","chapitre":"1SPE-EXPONENTIELLE","type_objet":"corrige","exercice_id":"1SPE-EXPO-EX-035","status":"generated"}
% BEGIN-VERIFY
% from sympy import *
% x = symbols('x')
% f = x*exp(x)
% fp = diff(f, x)
% assert simplify(fp - (x + 1)*exp(x)) == 0
% assert solve(fp, x) == [-1]
% assert f.subs(x, -1) == -exp(-1)
% END-VERIFY
\begin{corrige}{1SPE-EXPO-EX-035}

\textbf{Question 1.} On pose $u(x) = x$ et $v(x) = \mathrm{e}^{x}$.

On a $u'(x) = 1$ et $v'(x) = \mathrm{e}^{x}$.

Par la formule du produit :
\[f'(x) = u'(x)\,v(x) + u(x)\,v'(x) = 1 \times \mathrm{e}^{x} + x \times \mathrm{e}^{x}.\]

\textbf{Question 2.} On factorise par $\mathrm{e}^{x}$ :
\[f'(x) = \mathrm{e}^{x} + x\,\mathrm{e}^{x} = \boxed{(x+1)\,\mathrm{e}^{x}}.\]

\textbf{Question 3.} On résout $f'(x) = 0$, soit $(x+1)\,\mathrm{e}^{x} = 0$.

Comme $\mathrm{e}^{x} > 0$ pour tout réel $x$, le produit est nul si et seulement si $x + 1 = 0$, c'est-à-dire $\boxed{x = -1}$.

\textbf{Question 4.} Le signe de $f'(x)$ est celui de $x + 1$ (car $\mathrm{e}^{x} > 0$).
\begin{itemize}
\item Si $x < -1$, alors $f'(x) < 0$ : $f$ est décroissante.
\item Si $x > -1$, alors $f'(x) > 0$ : $f$ est croissante.
\end{itemize}
La fonction $f$ admet un minimum en $x = -1$, avec $f(-1) = -1 \times \mathrm{e}^{-1} = -\dfrac{1}{\mathrm{e}}$.
\end{corrige}
